\subsection{The arithmetic recursion with two higher-weight factors}

\begin{statement}{A.62}{Proved}{Ordinary ghosts with both higher-weight valuation factors}
For every positive integer $C$ divisible by $24$ there
is an infinite-support ordinary completed unit-equation
pair with the following explicit valuation pattern.
Choose $J\ge6$ at least as large as every prime
divisor of $C$, and put
\[
 n_j=(j+J)!,\quad D_j=6Cn_j^3,\quad
 E_j(Q)=Cn_j(Q^2-1)(Q^3-1).
 \tag{A.62.1}
\]
There are distinct primes $q_j,r_j,s_j>3$ for which
\[
 A=3\prod_{j\ge0}\frac{q_j^{E_j(q_j)}}{s_j^{E_j(s_j)}},
 \qquad
 B=-2\prod_{j\ge0}\frac{r_j^{E_j(r_j)}}{s_j^{E_j(s_j)}}.
 \tag{A.62.2}
\]
Every finite localization is the Kummer pair of an
ordinary $x_p\in\Q_p\setminus\{0,1\}$.
At two and three its valuation pairs are $(0,1)$
and $(1,0)$; at $q_j,r_j,s_j$ they are
$(E_j(q_j),0)$, $(0,E_j(r_j))$, and
$(-E_j(s_j),-E_j(s_j))$, respectively.
Thus this construction supplies both factors
$Q^2-1$ and $Q^3-1$ with a factorial factor
tending to zero in $\widehat\Z$.
The class-four construction uses $C=24$ and $J=6$.
\end{statement}
\begin{proof}
We prove the precise changes to the three-prime
construction A.59. Throughout, the coefficients
remain
$a=3u^{24}$, $b=-2v^{24}$, $c=w^{24}$
with pairwise disjoint supports away from two
and three. Since every exponent in (A.62.1)
is a multiple of $24$, these forms and the
three exact radical-field intersections
(A.59.4)--(A.59.6) remain valid without alteration.
We retain unit $D_j$-points at every finite
prime and high precision of $a+b-c$ at primes
up to $B(D_j)=(D_j+1)^4$.

\emph{The exact cubic local image.}
Fix $p\mid n$, write $s=v_p(n)$, and put
\[
 z_0=1+p^s,\quad
 \mathcal B=z_0+p^{s+1}\Z_p,\qquad
 F_n(z)=
 \exp\!\left(Cn(z^2-1)(z^3-1)\log z\right).
 \tag{A.62.3}
\]
Our starting factorial gives $s\ge4$ at two
and $s\ge2$ at three.
On positive integers in this ball, $F_n$ is
the ordinary power $z^{Cn(z^2-1)(z^3-1)}$.
Put
\[
 A_p=6Cn p^{3s},\qquad h=v_p(A_p)=4s+v_p(6C).
 \tag{A.62.4}
\]
There is an integral restricted power series
$R\in\Z_p\langle w\rangle$ such that
\[
 F_n(1+p^s w)=1+A_p f(w),
                       \qquad f(w)=w^3+p^sR(w).
 \tag{A.62.5}
\]
To verify this assertion, write $t=p^s w$.
The exponent in (A.62.3) is
$6Cn t^3V(t)$, where
\[
 V(t)=\left(1+\frac32t+\frac56t^2+\frac16t^3\right)
                         \frac{\log(1+t)}t.
 \tag{A.62.6}
\]
The linear coefficient of $V$ is one.
For $k\ge2$, the denominator valuation of
its $k$th coefficient is at most
$v_p(6)+\lfloor\log_p(k+1)\rfloor$.
Thus the coefficient after substitution
$t=p^s w$ has valuation at least $s$:
for $p\ge5$ use $s\ge1$, at three use
$s\ge2$, and at two use $s\ge4$.
These valuations tend to infinity.
Consequently $V(p^s w)=1+p^sR_0(w)$.
Expanding the exponential and dividing
by $A_p$, the further coefficients
$A_p^{m-1}/m!$ also belong to $p^s\Z_p$
and tend to zero. This proves (A.62.5).

The image of (A.62.3) on the required ball is exactly
\[
       F_n(\mathcal B)
           =F_n(z_0)+p^{h+1+v_p(3)}\Z_p.
 \tag{A.62.7}
\]
For $p\ne3$, substitute $w=1+pt$ into (A.62.5):
\[
       \frac{f(1+pt)-f(1)}p=3t+pR_1(t).
\]
The right side is a bijection of $\Z_p$
by the one-Lipschitz contraction argument.
At three one instead has
\[
       \frac{f(1+3t)-f(1)}9
              =t+3t^2+3t^3+3^{s-1}R_2(t),
                      \qquad R_2\in\Z_3\langle t\rangle.
\]
Here $s\ge2$ makes the last three terms
divisible by three, so this map too is a
bijection. These equations prove (A.62.7).
They also show
\[
                    v_p(F'_n(z_0))
                        =3s+v_p(18C).
 \tag{A.62.8}
\]
In particular the larger starting factorial
has supplied the precision needed at three;
no derivative dominance at $s=1$ is assumed.

At least one of $a,b,c$ is a unit at every prime.
The three-variable image on $\mathcal B^3$
is therefore
$(a+b-c)F_n(z_0)+p^{h+1+v_p(3)}\Z_p$,
by varying a variable with unit coefficient.
Define
\[
 d_j(p)=
 \begin{cases}
  3v_p(n_j)+v_p(18C),&p\mid n_j,\\
  0,&p\nmid n_j,
 \end{cases}
 \qquad H_j(p)=2d_j(p)+j+3.
 \tag{A.62.9}
\]
The inequality
$H_j(p)\ge h+1+v_p(3)$ follows by subtraction:
the difference is
$2s+v_p(18C)+j+2\ge0$.
Thus this precision of $a+b-c$ supplies the
local solution at every $p\mid n_j$.
Approximation to arbitrarily large later
CRT precision is allowed: (A.62.5) is
one-Lipschitz in $w$ after its displayed
factor, so a change of $z$ modulo $p^L$
changes $F_n(z)$ modulo $p^{h+L-s}$.

\emph{The prime constraints and their power images.}
At stage $j$ set $n=n_j$, and process every
prime $p\le P$, where $P$ is at least
$B(D_{j+1})$ and every old coefficient prime.
Choose the precisions as in A.59, at least
$H_{j+1}(p)$ and with $p^{L_p}>P$;
at $p\mid n$ also require $L_p\ge v_p(n)+2$.
The required residue conditions are now
\[
 \begin{array}{ll}
 v_p(z-1)=v_p(n),&p\mid n,\\
 z\ne0,\quad z^2\ne1,\quad z^3\ne1\pmod p,
                                             &p\nmid n.
 \end{array}
 \tag{A.62.10}
\]
Every prime in the second line is at least seven.
At a prime $\ell\mid n$, the elementary
lifting-the-exponent formulas give
\[
 v_\ell\bigl((z^2-1)(z^3-1)\bigr)
                     =2v_\ell(n)+v_\ell(6).
 \tag{A.62.11}
\]
For odd $\ell$ this follows by expanding
$z=1+\ell^{v_\ell(n)}u$ in the two factors;
the factor with exponent divisible by
$\ell$ gains exactly $v_\ell(6)$.
At two, $v_2(z+1)=1$ and
$z^2+z+1$ is odd, which gives the same formula.
At $\ell\nmid n$, (A.62.10) gives valuation zero.
Since every prime divisor of $6C$ already
divides $n$, these assertions give
\[
 v_\ell\bigl(Cn(z^2-1)(z^3-1)\bigr)
                         =v_\ell(6Cn^3).
 \tag{A.62.12}
\]
As in A.59, earlier choices determine each
exponent modulo $p-1$, because every prime
divisor of $p-1$ is smaller than $p$ and
its full required power is already controlled.
Hence the three exponent maps on
$\mathbf F_p^\times$ again have the common
image $(\mathbf F_p^\times)^{D_j}$.

At a small processed prime $p\nmid n$ choose
a generator $\lambda$ of this common cyclic
image. Every exponent map has a primitive
preimage of $\lambda$ in $\mathbf F_p^\times$.
For a proof, choose a generator $\zeta$ of
$\mathbf F_p^\times$, put $N=p-1$,
$g=\gcd(e,N)$ and $M=N/g$, and write
$\lambda=\zeta^{gt}$ with $(t,M)=1$.
The congruence
$(e/g)k\equiv t\pmod M$ has a unit solution
modulo $M$. Such a solution lifts to a unit
modulo $N$: at primes dividing $N$ but not
$M$ prescribe a nonzero residue by CRT;
at primes already dividing $M$ the residue
is a unit automatically.
Then $\zeta^k$ is the required primitive
preimage. This also covers the trivial
image group $M=1$.
Since $p-1>3$, a primitive preimage satisfies
both exclusions in (A.62.10).

On the residue ball of such a preimage $u$,
with Teichmuller lift $\omega$ and already
fixed exponent $e$, the local map is
\[
 {\cal F}_{p,e}(z)=\omega^e
       \exp\!\left(Cn(z^2-1)(z^3-1)
                                      \log(z/\omega)\right).
 \tag{A.62.13}
\]
With $z=\omega(1+pt)$, its logarithm divided
by $p$ is
\[
 Cn(\omega^2-1)(\omega^3-1)t+pG(t),
                            \qquad G\in\Z_p\langle t\rangle.
 \tag{A.62.14}
\]
The leading coefficient is a unit by
(A.62.10); thus the map is a bijection
onto $\omega^e(1+p\Z_p)$.
All three maps have a common Teichmuller
value lifting $\lambda$. The high
sum-error at small primes is canceled
by varying a variable with unit coefficient,
exactly as in (A.59.15)--(A.59.18).

At a processed prime $p>B(D_j)$, start
with the existing homogeneous unit
$D_j$-point. For each exponent map,
the set of images of the excluded residues
$u^2=1$ or $u^3=1$ has at most five
elements. Common scaling of the three
$D_j$-coordinates must therefore avoid
at most fifteen values in the common
power-image group.
That group has more than fifteen elements:
$(p-1)/\gcd(D_j,p-1)>15$ for
$p>(D_j+1)^4$ and $D_j\ge2$.
Choose the scaling outside the forbidden
set. All residue preimages of the scaled
values now satisfy (A.62.10).
The bijections (A.62.13)--(A.62.14)
lift them to the exact three ordinary
$D_j$-power values. This proves the
local equation at every processed large
prime, including old coefficient primes.

\emph{Primes, invariants, and limits.}
After these triangular choices, each
CRT class is a unit class equal to one
modulo $n$. The radical fields (A.59.4)
therefore supply distinct primes $q,r,s>P$,
using the old ratios $b/c,a/c,-a/b$.
Update by $a'=a q^{E_j(q)}$,
$b'=b r^{E_j(r)}$, $c'=c s^{E_j(s)}$.
The approximations give the target
$H_{j+1}(p)$ at every processed prime.
This bound exceeds $2v_p(D_{j+1})$,
so it also gives unit $D_{j+1}$-points
there by Hensel at $(1,1,1)$.

At each new coefficient prime $Q$,
\[
                      \gcd(D_{j+1},Q-1)\mid n.
 \tag{A.62.15}
\]
Indeed, prime divisors already in $n$
have exactly their old valuations in
$Q-1$. A new prime divisor of
$D_{j+1}$ comes from the next factorial
and is at most $j+J+1\le P$.
It cannot divide $Q-1$ by (A.62.10).
Every update exponent is divisible by
$n$, so the split old ratios remain
$n$th powers at the other new primes.
Their $n$th-power and $D_{j+1}$th-power
residue images coincide by (A.62.15).
The three simple-root Hensel arguments
following (A.59.21) consequently give
unit $D_{j+1}$-points at all three new
coefficient primes.

At every unprocessed remaining prime,
the coefficients are units and the
prime exceeds $B(D_{j+1})$.
The Jacobi-sum bound in A.44 gives
unit-coordinate $D_{j+1}$-points.
This completes the all-place invariant
and the high-error invariant, starting
from $(3,-2,1)$ and $(1,1,1)$.

Finally, the exponents are all multiples
of the cofinal factorials $n_j$.
The completed products converge, and
at each fixed finite prime the three
integer coefficient products have
eventually fixed valuations and
convergent unit parts. Their ordinary
limits satisfy $a_p+b_p=c_p\ne0$
by the growing error precision.
Thus $x_p=a_p/c_p$ and $1-x_p=b_p/c_p$
give precisely the localizations of
(A.62.2). The signs at the real place
are again represented by $(3,-2)$.
The disjoint new prime supports give
exactly the stated valuation pairs
and their infinite support.
All these limit arguments use the
same completed and local maps as
(A.59.22)--(A.59.23).
\end{proof}

The exact certificate
\texttt{code/cubic\_exponent\_certificate.py}
checks the exponent valuations and triangular
gcd identities, primitive preimages of image
generators, the full cubic local image on
finite quotient balls by literal modular
integer powers, the restricted-series
coefficient bounds, and all stated precision
inequalities. Its output is retained.
No complete large prime-search stage was
executed; its termination is supplied by
the proved field intersections and Chebotarev.
